% =====================================================================
\chapter{Polinomios, ecuaciones y sistemas}
\label{tema:3}
% =====================================================================

Los polinomios son las funciones más sencillas después de las lineales y aparecen por todas partes: como aproximaciones de Taylor, como funciones de pérdida cuadráticas y como polinomio característico de una matriz. Este tema repasa su aritmética (división, raíces, factorización), la resolución de ecuaciones e inecuaciones y los sistemas lineales pequeños. Sigue los capítulos~2, 5 y~7 de OpenStax \emph{College Algebra}~\cite{openstax}, con las demostraciones que allí se omiten.

\section{Polinomios y división euclídea}

\begin{definicion}[Polinomio]
Un \emph{polinomio} con coeficientes reales es una expresión $P(x)=a_nx^n+\dots+a_1x+a_0$ con $a_i\in\R$. Si $a_n\neq0$, $n$ es el \emph{grado} de $P$, $\deg P$, y $a_n$ su \emph{coeficiente principal}. Un número $r$ es una \emph{raíz} de $P$ si $P(r)=0$.
\end{definicion}

Si $P$ y $Q$ son no nulos, $\deg(PQ)=\deg P+\deg Q$, porque el coeficiente principal del producto es el producto de los coeficientes principales, que no es cero (\cref{prop:cuerpo}(iv)).

\begin{teorema}[División euclídea]\label{teo:division}
Sean $P$ y $D$ polinomios con $D\neq0$. Existen polinomios únicos $Q$ (cociente) y $R$ (resto) tales que
\[
P=DQ+R\qquad\text{y}\qquad R=0\ \text{o}\ \deg R<\deg D.
\]
\end{teorema}

\begin{proof}
\emph{Existencia}, por inducción sobre $\deg P$. Si $P=0$ o $\deg P<\deg D$, se toma $Q=0$ y $R=P$. Si $\deg P=n\ge m=\deg D$, con coeficientes principales $a_n$ y $d_m$, el polinomio $P_1=P-\frac{a_n}{d_m}x^{n-m}D$ tiene grado menor que $n$ (o es nulo), porque se ha cancelado el término de grado $n$. Por la hipótesis de inducción, $P_1=DQ_1+R$ con $R$ como se pide, y entonces $P=D\bigl(\frac{a_n}{d_m}x^{n-m}+Q_1\bigr)+R$.

\emph{Unicidad}. Si $DQ+R=DQ'+R'$, entonces $D(Q-Q')=R'-R$. Si $Q\neq Q'$, el miembro izquierdo tiene grado al menos $\deg D$ y el derecho grado menor que $\deg D$, lo cual es imposible. Luego $Q=Q'$ y, en consecuencia, $R=R'$.
\end{proof}

\begin{teorema}[Teoremas del resto y del factor]\label{teo:resto}
El resto de dividir $P$ entre $x-a$ es el número $P(a)$. En consecuencia, $a$ es raíz de $P$ si y solo si $x-a$ divide a $P$.
\end{teorema}

\begin{proof}
Por el \cref{teo:division}, $P(x)=(x-a)Q(x)+R$ con $R$ constante, porque $\deg R<1$. Evaluando en $x=a$, $P(a)=0\cdot Q(a)+R=R$.
\end{proof}

La \emph{regla de Ruffini} es una forma abreviada de dividir entre $x-a$: se bajan los coeficientes uno a uno, multiplicando el último resultado por $a$ y sumándolo al coeficiente siguiente.

\begin{ejemplo}\label{ej:ruffini}
Dividir $P(x)=2x^3-5x^2+0x+7$ entre $x-2$:
\[
\begin{array}{r|rrrr}
 & 2 & -5 & 0 & 7\\
2 & & 4 & -2 & -4\\\hline
 & 2 & -1 & -2 & 3
\end{array}
\]
Cociente $2x^2-x-2$ y resto $3$; en efecto, $P(2)=16-20+7=3$. El coeficiente nulo de $x$ debe escribirse.
\end{ejemplo}

\begin{corolario}\label{cor:raices}
Un polinomio no nulo de grado $n$ tiene a lo sumo $n$ raíces reales distintas.
\end{corolario}

\begin{proof}
Por inducción sobre $n$. Un polinomio de grado~0 es una constante no nula y no tiene raíces. Si $\deg P=n\ge1$ y $P$ tiene alguna raíz $a$, por el \cref{teo:resto} $P=(x-a)Q$ con $\deg Q=n-1$. Si $b\neq a$ es otra raíz, $0=(b-a)Q(b)$ obliga a $Q(b)=0$. Así, las raíces de $P$ son $a$ y las de $Q$, que por la hipótesis de inducción son a lo sumo $n-1$.
\end{proof}

\begin{teorema}[Raíces racionales]\label{teo:raices-racionales}
Sea $P(x)=a_nx^n+\dots+a_0$ con coeficientes enteros, $a_n\neq0$ y $a_0\neq0$. Si $p/q$, con $p,q$ enteros primos entre sí, es raíz de $P$, entonces $p$ divide a $a_0$ y $q$ divide a $a_n$.
\end{teorema}

\begin{proof}
Multiplicando $P(p/q)=0$ por $q^n$: $a_np^n+a_{n-1}p^{n-1}q+\dots+a_1pq^{n-1}+a_0q^n=0$. Todos los términos salvo el último son múltiplos de $p$, luego $p\mid a_0q^n$; como $p$ y $q$ no tienen factores primos comunes, $p\mid a_0$. Del mismo modo, todos los términos salvo el primero son múltiplos de $q$, luego $q\mid a_np^n$ y $q\mid a_n$.
\end{proof}

\section{La ecuación de segundo grado y las fórmulas de Vieta}

\begin{teorema}[Ecuación de segundo grado]\label{teo:cuadratica}
Sean $a\neq0$ y $\Delta=b^2-4ac$ el \emph{discriminante}. La ecuación $ax^2+bx+c=0$ tiene dos soluciones reales distintas si $\Delta>0$, una (doble) si $\Delta=0$ y ninguna si $\Delta<0$. Cuando $\Delta\ge0$,
\[
x=\frac{-b\pm\sqrt{\Delta}}{2a}.
\]
\end{teorema}

\begin{proof}
Se completa el cuadrado:
\[
ax^2+bx+c=a\Bigl[\Bigl(x+\frac{b}{2a}\Bigr)^2-\frac{\Delta}{4a^2}\Bigr].
\]
Como $a\neq0$, la ecuación equivale a $(x+\frac b{2a})^2=\frac{\Delta}{4a^2}$. El cuadrado de un real es no negativo, luego no hay solución si $\Delta<0$; si $\Delta\ge0$, extrayendo la raíz se obtiene $x+\frac b{2a}=\pm\frac{\sqrt\Delta}{2a}$.
\end{proof}

\begin{proposicion}[Fórmulas de Vieta]\label{prop:vieta}
Si $x_1,x_2$ son las raíces (reales o repetidas) de $ax^2+bx+c$, entonces
\[
ax^2+bx+c=a(x-x_1)(x-x_2),\qquad x_1+x_2=-\frac ba,\qquad x_1x_2=\frac ca.
\]
Más en general, si $P(x)=a_n\prod_{i=1}^n(x-x_i)$, entonces $\sum_ix_i=-a_{n-1}/a_n$ y $\prod_ix_i=(-1)^na_0/a_n$.
\end{proposicion}

\begin{proof}
Si $\Delta\ge0$, la expresión obtenida al completar el cuadrado en la demostración del \cref{teo:cuadratica} es una diferencia de cuadrados:
\[
ax^2+bx+c=a\Bigl(x+\frac b{2a}-\frac{\sqrt\Delta}{2a}\Bigr)\Bigl(x+\frac b{2a}+\frac{\sqrt\Delta}{2a}\Bigr)=a(x-x_1)(x-x_2).
\]
Desarrollando, $a(x-x_1)(x-x_2)=ax^2-a(x_1+x_2)x+ax_1x_2$, e igualando coeficientes se obtienen las fórmulas. El caso general es análogo: en $a_n\prod_i(x-x_i)$, el coeficiente de $x^{n-1}$ es $-a_n\sum_ix_i$ y el término independiente es $a_n\prod_i(-x_i)$.
\end{proof}

\begin{ejemplo}
Para $x^2-7x+10$: $x_1+x_2=7$ y $x_1x_2=10$, lo que sugiere $x_1=2$, $x_2=5$; en efecto, $x^2-7x+10=(x-2)(x-5)$.
\end{ejemplo}

\section{Otras ecuaciones e inecuaciones}

Las ecuaciones \emph{bicuadradas} $ax^4+bx^2+c=0$ se reducen a segundo grado con $t=x^2$. En las ecuaciones \emph{con radicales} se aísla la raíz y se eleva al cuadrado; como elevar al cuadrado no es una operación reversible ($a^2=b^2$ no implica $a=b$), pueden aparecer soluciones \emph{extrañas} y es obligatorio comprobar cada candidata en la ecuación original. En las ecuaciones \emph{racionales} se multiplica por el denominador común y se descartan los valores que lo anulan.

Para resolver una inecuación $P(x)/Q(x)\ge0$ se factorizan numerador y denominador, se marcan en la recta los puntos donde se anula algún factor y se estudia el signo de cada factor en cada intervalo (\emph{tabla de signos}); los ceros del denominador nunca forman parte de la solución.

\begin{ejemplo}\label{ej:tabla-signos}
Resolver $\dfrac{x-1}{x+2}<0$. Los puntos críticos son $-2$ y $1$. En $(-\infty,-2)$ ambos factores son negativos y el cociente es positivo; en $(-2,1)$ el numerador es negativo y el denominador positivo; en $(1,\infty)$ ambos son positivos. La solución es $(-2,1)$.
\end{ejemplo}

\section{Sistemas lineales pequeños}

Un sistema de $m$ ecuaciones lineales con $n$ incógnitas es \emph{compatible determinado} si tiene una única solución, \emph{compatible indeterminado} si tiene infinitas e \emph{incompatible} si no tiene ninguna. Los métodos de sustitución, igualación y reducción de bachillerato son casos particulares de la eliminación gaussiana del módulo~3.

\begin{teorema}[Regla de Cramer para $2\times2$]\label{teo:cramer}
El sistema $ax+by=e$, $cx+dy=f$ tiene solución única si y solo si $\delta=ad-bc\neq0$, y entonces
\[
x=\frac{ed-bf}{\delta},\qquad y=\frac{af-ec}{\delta}.
\]
\end{teorema}

\begin{proof}
Multiplicando la primera ecuación por $d$, la segunda por $b$ y restando: $(ad-bc)x=ed-bf$. Multiplicando la primera por $c$, la segunda por $a$ y restando: $(ad-bc)y=af-ec$. Si $\delta\neq0$, estas igualdades determinan $x$ e $y$, y se comprueba por sustitución que son solución. El caso $\delta=0$ se estudia en el \cref{eje:3-det-cero}.
\end{proof}

\begin{ejemplo}\label{ej:gauss3}
El sistema $x+y+z=6$, $2x-y+z=3$, $x+2y-z=2$ se resuelve por reducción. Restando a la segunda el doble de la primera y a la tercera la primera: $-3y-z=-9$ y $y-2z=-4$. De la última, $y=2z-4$; sustituyendo, $-3(2z-4)-z=-9$, es decir, $-7z=-21$ y $z=3$. Entonces $y=2$ y $x=1$.
\end{ejemplo}

\begin{observacion}[Errores frecuentes]
\begin{enumerate}[label=(\alph*)]
\item Dividir ambos miembros por $x$ o por $(x-a)$ pierde la solución $x=0$ o $x=a$; es preferible factorizar.
\item En las ecuaciones con radicales, no comprobar las soluciones produce soluciones extrañas (\cref{eje:3-ecuaciones}).
\item Multiplicar una inecuación por el denominador sin conocer su signo invierte, o no, el sentido; la tabla de signos evita el problema.
\item En la regla de Ruffini hay que escribir los coeficientes nulos.
\end{enumerate}
\end{observacion}

\begin{observacion}[Conexión con el curso]
Los autovalores de una matriz $2\times2$ son las raíces de su polinomio característico $\lambda^2-(\operatorname{tr}A)\lambda+\det A$; por las fórmulas de Vieta, su suma es la traza y su producto el determinante (módulo~3). Completar el cuadrado es la forma más rápida de minimizar una pérdida cuadrática en una variable, que es el modelo de juguete con el que se estudia la tasa de aprendizaje en el módulo~7. Las ecuaciones normales de la regresión lineal (módulo~8) son un sistema lineal.
\end{observacion}

\begin{lecturas}
\item OpenStax \emph{College Algebra}~\cite{openstax}, §1.4--1.5 (polinomios y factorización), §2.5 «Quadratic Equations», §2.6 «Other Types of Equations».
\item OpenStax~\cite{openstax}, §5.4 «Dividing Polynomials» y §5.5 «Zeros of Polynomial Functions» (teoremas del resto, del factor y de las raíces racionales).
\item OpenStax~\cite{openstax}, §7.1--7.2 (sistemas lineales de dos y tres incógnitas) y §7.8 «Solving Systems with Cramer's Rule».
\end{lecturas}

\section*{Ejercicios del tema 3}
\addcontentsline{toc}{section}{Ejercicios del tema 3}

\begin{ejercicio}{1}{eje:3-basico}
Sea $P(x)=2x^3-x+5$. (a) Calcule $P(-1)$. (b) ¿Qué grado tiene $P(x)^2\cdot(x^2+1)$? (c) Escriba una ecuación de segundo grado con coeficientes enteros cuyas raíces sean $3$ y $-\frac12$.
\end{ejercicio}

\begin{ejercicio}{1}{eje:3-discriminante}
Sin resolverlas, determine cuántas soluciones reales tiene cada ecuación: (a) $x^2-4x+5=0$; (b) $4x^2-12x+9=0$; (c) $2x^2+x-3=0$.
\end{ejercicio}

\begin{ejercicio}{2}{eje:3-division}
Divida $P(x)=2x^4-3x^3+x-5$ entre $D(x)=x^2-x+1$ y compruebe el resultado.
\fuente{adaptado de OpenStax \emph{College Algebra}~\cite{openstax}, §5.4.}
\end{ejercicio}

\begin{ejercicio}{2}{eje:3-ruffini}
Factorice con la regla de Ruffini: (a) $x^3-6x^2+11x-6$; (b) $x^4-5x^2+4$.
\fuente{adaptado de OpenStax \emph{College Algebra}~\cite{openstax}, §5.4--5.5.}
\end{ejercicio}

\begin{ejercicio}{2}{eje:3-resto}
(a) Halle $k$ para que $x^3+kx^2-4x+4$ sea divisible por $x-2$. (b) Halle, sin dividir, el resto de $x^{100}-1$ entre $x+1$ y el de $x^{101}+3$ entre $x+1$.
\end{ejercicio}

\begin{ejercicio}{2}{eje:3-racionales}
Halle todas las raíces de $2x^3-3x^2-3x+2$ usando el \cref{teo:raices-racionales}.
\fuente{adaptado de OpenStax \emph{College Algebra}~\cite{openstax}, §5.5 («Rational Zero Theorem»).}
\end{ejercicio}

\begin{ejercicio}{2}{eje:3-vieta}
Sean $x_1,x_2$ las raíces de $x^2-5x+3=0$. Sin calcularlas, halle $x_1^2+x_2^2$, $\dfrac1{x_1}+\dfrac1{x_2}$ y $x_1^3+x_2^3$.
\end{ejercicio}

\begin{ejercicio}{2}{eje:3-ecuaciones}
Resuelva: (a) $x^4-13x^2+36=0$; (b) $\sqrt{x+5}=x-1$; (c) $\dfrac1{x-1}+\dfrac2{x+1}=1$.
\fuente{adaptado de OpenStax \emph{College Algebra}~\cite{openstax}, §2.6.}
\end{ejercicio}

\begin{ejercicio}{2}{eje:3-inecuaciones}
Resuelva: (a) $x^2-x-6<0$; (b) $\dfrac{(x-1)(x+2)}{x-3}\ge0$.
\end{ejercicio}

\begin{ejercicio}{2}{eje:3-sistemas}
(a) Discuta y resuelva, según el valor de $k$, el sistema $x+2y=4$, $3x+ky=6$. (b) Resuelva $x+y+z=6$, $2x-y+z=3$, $x+2y-z=2$ por un orden de eliminación distinto del del \cref{ej:gauss3}.
\fuente{adaptado de OpenStax \emph{College Algebra}~\cite{openstax}, §7.1--7.2 y §7.8.}
\end{ejercicio}

\begin{ejercicio}{3}{eje:3-identidad}
Demuestre que si dos polinomios $P$ y $Q$ de grado a lo sumo $n$ toman el mismo valor en $n+1$ puntos distintos, entonces $P=Q$ (tienen los mismos coeficientes).
\end{ejercicio}

\begin{ejercicio}{3}{eje:3-diferencia-potencias}
Demuestre que, para todo $n\in\N$,
\[
x^n-a^n=(x-a)(x^{n-1}+ax^{n-2}+\dots+a^{n-2}x+a^{n-1}).
\]
Deduzca que $x-a$ divide a $x^n-a^n$ y que, si $n$ es impar, $x+a$ divide a $x^n+a^n$.
\end{ejercicio}

\begin{ejercicio}{3}{eje:3-impares}
Sean $b$ y $c$ enteros impares. Demuestre que la ecuación $x^2+bx+c=0$ no tiene soluciones racionales.
\end{ejercicio}

\begin{ejercicio}{3}{eje:3-det-cero}
Demuestre que si $ad-bc=0$, el sistema $ax+by=e$, $cx+dy=f$ o bien no tiene solución o bien tiene infinitas. (Suponga que algún coeficiente $a,b,c,d$ es no nulo.)
\end{ejercicio}

\begin{ejercicio}{3}{eje:3-cuadratica-positiva}
Sea $a>0$ y $f(x)=ax^2+bx+c$. Demuestre que $f(x)\ge c-\dfrac{b^2}{4a}$ para todo $x$, con igualdad solo en $x=-\dfrac{b}{2a}$, y deduzca que $f(x)\ge0$ para todo $x\in\R$ si y solo si $b^2-4ac\le0$.
\end{ejercicio}
